\section{Weyl semimetals}

The Weyl semimetals have nondegenerate gapless linear dispersions around band touching points (Weyl nodes). The low-energy effective Hamiltonian around a Weyl node is written as:
\begin{equation}
H_W(\bm{k}) = Q v_F \bm{k} \cdot \bm{\sigma}
\end{equation}
where $Q=\pm 1$ indicates the chirality, $v_F$ is the Fermi velocity and $\sigma_i$ are Pauli matrices. Because the sum of chiralities of Weyl nodes in a system must be zero, the simplest realization of a Weyl semimetal is one with two Weyl nodes of opposite chiralities. Note that the minimal number of Weyl nodes in Weyl semimetal with broken inversion symmetry is four, while it is two in Weyl semimetals with broken time-reversal symmetry.

Now, let's consider a $4\times 4$ continuum model Hamiltonian for two-node Weyl semimetals with broken time-reversal symmetry:
\begin{equation}
H_0(\bm{k}) = v_F (\tau_z \otimes \bm{k} \cdot \bm{\sigma} + \Delta \tau_x \otimes \mathbb{I} + \mathbb{I} \otimes \bm{b} \cdot \bm{\sigma})
\end{equation}
where $\tau_i$ and $\sigma_i$ are Pauli matrices for Weyl-node and spin degrees of freedom respectively and $\Delta$ is the mass of 3D Dirac fermions. The term $\bm{b} \cdot \bm{\sigma}$ represents a magnetic interaction such as exchange interaction between conduction electrons and magnetic impurities or the Zeeman coupling with an external magnetic field. Note the Hamiltonian with $\bm{b}=0$ describes a topological or normal insulator depending on the sign of $\Delta$ (mass). Without loss of generality, we may set $\bm{b}=(0,0,b)$.

Now let's find the location of Weyl nodes. The Hamiltonian of Weyl nodes can be written as:
\begin{equation}
\begin{aligned}
H &= v_F \tau_z (k_x \sigma_x + k_y \sigma_y + k_z \sigma_z) + \Delta \tau_x \otimes \mathbb{I} + b \mathbb{I} \otimes \sigma_z \\
&= \underbrace{v_F \tau_z \otimes (k_x \sigma_x + k_y \sigma_y)}_{H_\perp} + \underbrace{v_F \tau_z \otimes k_z \sigma_z + \Delta \tau_x \otimes \mathbb{I} + b \mathbb{I} \otimes \sigma_z}_{H_\parallel}
\end{aligned}
\end{equation}
Because $H_\perp$ and $H_\parallel$ anticommute, so we can write down:
\begin{equation}
\begin{aligned}
H^2 &= H_\perp^2 + H_\parallel^2 + \{H_\perp, H_\parallel\} \\
&= H_\perp^2 + H_\parallel^2
\end{aligned}
\end{equation}
This leads to the explicit energy spectrum:
\begin{equation}
E^2 = v_F^2 (k_x^2 + k_y^2) + E_\parallel^2
\end{equation}
and $H_\parallel$ can be split into two $2 \times 2$ blocks:
\begin{equation}
\begin{aligned}
h_{eff} &= s(k_z \tau_z + b) + \Delta \tau_x \\
&= s k_z \cdot \tau_z + \Delta \tau_x + s b \mathbb{I}
\end{aligned}
\end{equation}
where $s=\pm 1$. So the energy spectrum of $h_{eff}$ is given by:
\begin{equation}
E_\parallel = s b \pm \sqrt{k_z^2 + \Delta^2}
\end{equation}
So the Weyl node location satisfies:
\begin{equation}
    \left\{
    \begin{aligned}
    & k_x^2 + k_y^2 = 0 \\
    & s b \pm \sqrt{v_F^2 k_z^2 + \Delta^2} = 0
    \end{aligned}
    \right.
    \end{equation}
Thus, the coordinates of the Weyl nodes are:
\begin{equation}
(k_x, k_y, k_z) = (0, 0, \pm \frac{\sqrt{b^2-\Delta^2}}{v_F})
\end{equation}
which requires $b^2 > \Delta^2$.

And then we derive the $\theta$ term of Weyl semimetals from the microscopic four-band model. In order to describe a more generic Weyl semimetal, we add the term $\mu_5 \tau_z$ to the Hamiltonian which generates a chemical potential difference $2\mu_5$ between the two Weyl nodes. We also set $\Delta=0$ for simplicity. So the Hamiltonian can be written as:
\begin{equation}
H_0' = v_F (\tau_z \otimes \bm{\sigma} \cdot \bm{k} + \mathbb{I} \otimes \bm{\sigma} \cdot \bm{b}) - \mu_5 \tau_z \otimes \mathbb{I}
\end{equation}
We rewrite the Hamiltonian with Gamma matrices.

\begin{equation}
\left\{
\begin{aligned}
& \gamma^0 = \tau_x \otimes \mathbb{I} \\
& \gamma^i = -\mathrm{i} \tau_y \otimes \sigma_i \\
& \gamma^5 = \tau_z \otimes \mathbb{I} \\
& b_\mu = (\mu_5, -\bm{b})
\end{aligned}
\right.
\end{equation}
and we set $v_F \approx c = 1$, so:
\begin{equation}
H = - \gamma^0 \gamma^i k_i + \gamma^0 \gamma^i \gamma^5 b_i - \mu_5 \gamma^5
\end{equation}
The action of the system in the presence of external electric and magnetic fields with the four potential $A_\mu = (A_0, -\bm{A})$ is given by
\begin{equation}
\begin{aligned}
S &= \int \mathrm{d}t \mathrm{d}^3x \, \psi^\dagger [\mathrm{i}(\partial_t - \mathrm{i} e A_0) - H_0'(\bm{k}+e\bm{A})] \\
&= \int \mathrm{d}t \mathrm{d}^3x \, \bar{\psi} \gamma^0 [\mathrm{i} \partial_t + \mathrm{i} \gamma^0 \gamma^i \partial_i + e A_0 - e \gamma^0 \gamma^i A_i - \gamma^0 \gamma^i \gamma^5 b_i + \mu_5 \gamma^5] \psi \\
&= \int \mathrm{d}t \mathrm{d}^3x \, \bar{\psi} \mathrm{i} [\gamma^0 \partial_0 + \gamma^i \partial_i - \mathrm{i} e (\gamma^0 A_0 - \gamma^i A_i) - \mathrm{i} (\gamma^0 \mu_5 - \gamma^i b_i) \gamma^5] \psi \\
&= \int \mathrm{d}t \mathrm{d}^3x \, \bar{\psi} \mathrm{i} (\gamma^\mu \partial_\mu - \mathrm{i} e \gamma^\mu A_\mu - \mathrm{i} \gamma^\mu b_\mu \gamma^5) \psi \\
&= \int \mathrm{d}t \mathrm{d}^3x \, \bar{\psi} \mathrm{i} \gamma^\mu (\partial_\mu - \mathrm{i} e A_\mu - \mathrm{i} b_\mu \gamma^5) \psi
\end{aligned}
\end{equation}
where $e>0$, $\psi$ is a four-component spinor, $\bar{\psi} = \psi^\dagger \gamma^0$.

Now we apply Fujikawa's method to the action. Performing an infinitesimal gauge transformation for infinite times such that:
\begin{equation}
\left\{
\begin{aligned}
& \psi \to \psi' = \mathrm{e}^{-\frac{\mathrm{i}\theta(\bm{r},t)\gamma^5 \mathrm{d}\phi}{2}} \psi \\
& \bar{\psi} \to \bar{\psi}' = \bar{\psi} \mathrm{e}^{-\frac{\mathrm{i}\theta(\bm{r},t)\gamma^5 \mathrm{d}\phi}{2}}
\end{aligned}
\right.
\end{equation}
with $\theta(\bm{r},t) = 2 x^\mu b_\mu = 2(\bm{b} \cdot \bm{r} - \mu_5 t)$ and $\phi \in [0, 1]$, then the partition function $Z$ is transformed as:
\begin{equation}
\begin{aligned}
Z &= \int D[\psi, \bar{\psi}] \mathrm{e}^{\mathrm{i}S[\psi, \bar{\psi}]} \\
\Rightarrow Z' &= \int D[\psi', \bar{\psi}'] \mathrm{e}^{\mathrm{i}S[\psi', \bar{\psi}']}
\end{aligned}
\end{equation}
Because the anticommutation relation $\{\gamma^\mu, \gamma^5\} = 0$ is satisfied, so:
\begin{equation}
\mathrm{e}^{\frac{-\mathrm{i}\theta(\bm{r},t)\gamma^5 \mathrm{d}\phi}{2}} \gamma^\mu = \gamma^\mu \mathrm{e}^{\frac{\mathrm{i}\theta(\bm{r},t)\gamma^5 \mathrm{d}\phi}{2}}
\end{equation}
The Lagrangian of $S[\psi', \bar{\psi}']$ can be written as:
\begin{equation}
\begin{aligned}
\mathcal{L} &= \bar{\psi} \mathrm{e}^{-\frac{\mathrm{i}\theta(\bm{r},t)\gamma^5 \mathrm{d}\phi}{2}} \cdot \mathrm{i} \gamma^\mu (\partial_\mu - \mathrm{i} e A_\mu - \mathrm{i} b_\mu \gamma^5) \mathrm{e}^{-\frac{\mathrm{i}\theta(\bm{r},t)\gamma^5 \mathrm{d}\phi}{2}} \psi \\
&= \bar{\psi} \mathrm{i} \gamma^\mu (\partial_\mu - \mathrm{i} e A_\mu) \psi + \bar{\psi} \mathrm{i} \gamma^\mu (\partial_\mu \mathrm{e}^{-\frac{\mathrm{i}\theta(\bm{r},t)\gamma^5 \mathrm{d}\phi}{2}}) \psi - \bar{\psi} \mathrm{i} \gamma^\mu (-\mathrm{i} b_\mu \gamma^5) \psi \\
&= \bar{\psi} \mathrm{i} \gamma^\mu (\partial_\mu - \mathrm{i} e A_\mu) \psi - \bar{\psi} \mathrm{i} \gamma^\mu (1 - \mathrm{d}\phi) \mathrm{i} b_\mu \gamma^5 \psi
\end{aligned}
\end{equation}
And the Jacobian $J$ of transformation of $D[\psi', \bar{\psi}'] = J D[\psi, \bar{\psi}]$ is given by:
\begin{equation}
\ln J = \mathrm{i} \mathrm{d}\phi \int \mathrm{d}t \mathrm{d}^3x \frac{e^2}{16\pi^2} \cdot \theta(\bm{r}, t) \cdot {}^*F^{\mu\nu} F_{\mu\nu}
\end{equation}
We integrate with respect to the variable $\phi$ from 0 to 1. Due to the invariance of the partition function, we arrive at the following expression of $S$:
\begin{equation}
S = \int \mathrm{d}t \mathrm{d}^3x \, \bar{\psi} \mathrm{i} \gamma^\mu (\partial_\mu - \mathrm{i} e A_\mu) \psi + \frac{e^2}{16\pi^2} \int \mathrm{d}t \mathrm{d}^3x \, \theta(\bm{r},t) {}^*F^{\mu\nu} F_{\mu\nu}
\end{equation}
where the first term represents the trivial action of massless Dirac fermions and the second term is the $\theta$ term with $\theta(\bm{r},t) = 2(\bm{b} \cdot \bm{r} - \mu_5 t)$.

\subsection{Anomalous Hall effect and negative magnetoresistance}

Let's consider the consequences of the presence of a $\theta$ term in Weyl semimetals. As we have also seen in the case of insulators with a dynamical axion field, an electric current is induced in the presence of a $\theta$ term. The induced electric current density and charge density are given by:
\begin{equation}
\left\{
\begin{aligned}
& \bm{j}(\bm{r}, t) = -\frac{\alpha}{\pi} (\dot{\theta} \cdot \bm{B} + \nabla \theta \times \bm{E}) \\
& \rho(\bm{r}, t) = \frac{\alpha}{\pi} (\nabla \theta) \cdot \bm{B}
\end{aligned}
\right.
\end{equation}
    
In the present case of $\theta(\bm{r}, t) = 2(\bm{b} \cdot \bm{r} - \mu_5 t)$, we readily obtain a static current of the form:
\begin{equation}
\bm{j} = \frac{2\alpha}{\pi} (\mu_5 \bm{B} - \bm{b} \times \bm{E})
\end{equation}
The second term is the anomalous Hall effect. Let's consider a 2D plane which is perpendicular to $\bm{b}$. And we set $\bm{b}=(0,0,b)$, which gives the 3D Hall conductivity as:
\begin{equation}
\sigma_{xy}^{3D} = \frac{2\alpha b}{\pi}
\end{equation}
The first term is the chiral magnetic effect. The chiral magnetic effect indicates that a direct current is generated along a static magnetic field even in the absence of electric fields, when there exists a chemical potential difference $\delta\mu=2\mu_5$ between the two Weyl nodes. However, the existence of the static chiral magnetic effect is ruled out in crystalline solid. But the chiral magnetic effect can be realized under nonequilibrium circumstance when the system is driven from equilibrium by the combined effect of electric and magnetic fields. In this case of Weyl semimetal potential difference can be generated by the chiral anomaly. The chiral anomaly in weyl semimetals is referred to as the electron number nonconservation in a given weyl cone under parallel electric and magnetic fields in which the rate of pumping of electrons is given by:
\begin{equation}
\frac{\partial N_i}{\partial t} = Q_i \frac{\alpha}{\pi} \bm{E} \cdot \bm{B}
\end{equation}
where $i$ is a valley (Weyl node) index and
\begin{equation}
Q_i = \int \frac{\mathrm{d}^3k}{2\pi} \cdot \frac{\partial f_0(\mathcal{E}_{\bm{k}}^m)}{\partial \mathcal{E}_{\bm{k}}^m} \cdot \bm{v}_{\bm{k}}^m \cdot \bm{\Omega}_{\bm{k}}^m
\end{equation}
is the chirality of the valley. Here, $\mathcal{E}_{\bm{k}}^m$ is the energy of Bloch electrons with momentum $\bm{k}$ in band $m$ in a given valley $i$, $f_0(\mathcal{E}_{\bm{k}}^m)$ is the Fermi distribution function. $\bm{v}_{\bm{k}}^m$ is the group velocity and $\bm{\Omega}_{\bm{k}}^m$ is the Berry curvature. So the difference of the total electron number between the weyl nodes leads to the difference of the chemical potential between weyl nodes.

A phenomenon manifested by the chiral anomaly is a negative magnetoresistance:
\begin{equation}
\sigma_{zz} = \frac{\alpha}{\pi} \frac{(e B_z)^2 v_F^3}{\mu^2} \tau_{\mathrm{inter}}
\end{equation}
where $\mu$ is the equilibrium chemical potential and $\tau_{\mathrm{inter}}$ is the intervalley scattering time. This unusual magnetoconductivity holds in the low-field limit $B_z \to 0$. So the expressing of $\bm{j}$ is understood as coming from
\begin{equation}
\bm{j} \propto (\bm{E} \cdot \bm{B} \tau_{\mathrm{inter}}) \cdot \bm{B}
\end{equation}
which indicates that it is a consequence of the chiral magnetic effect.